\begin{table}[htbp]
    \centering
    \caption{两个毕业设计候选选题对比}
    \label{tab:topic_comparison}
    \renewcommand{\arraystretch}{1.4}
    \setlength{\tabcolsep}{6pt}

    \begin{tabularx}{\textwidth}{
        >{\centering\arraybackslash}p{2.4cm}
        >{\raggedright\arraybackslash}X
        >{\raggedright\arraybackslash}X
    }

        \hline
        \textbf{对比维度}
        & \textbf{选题一：鄱阳湖候鸟目标检测}
        & \textbf{选题二：无人机红外人员搜救} \\
        \hline

        应用场景
        & 鄱阳湖湿地候鸟生态监测
        & 无人机灾害环境人员搜救 \\

        检测对象
        & 白鹤、白额雁、白鹳、普通鹤、天鹅等典型候鸟
        & 复杂环境下的人员目标 \\

        数据集
        & PYL-5-2023鄱阳湖鸟类数据集
        & POP红外人员检测数据集 \\

        数据规模
        & 8077张高质量鸟类图像
        & 8768张红外图像、25811个目标标注框 \\

        数据标注
        & X-AnyLabeling人工标注，包含类别和位置信息
        & COCO格式标注，并提供格式转换代码 \\

        数据来源
        & 来源于鄱阳湖湿地生态系统监测预警平台历史监控录像
        & 基于无人机红外成像采集，数据集公开发布 \\

        图像类型
        & 可见光监控图像
        & 红外热成像图像 \\

        主要技术问题
        & 小目标、姿态变化、遮挡、复杂背景、多目标检测
        & 小目标、植被遮挡、复杂背景、不同飞行高度 \\

        模型研究
        & 基于轻量化YOLO模型，针对鸟类小目标检测进行优化
        & 基于轻量化YOLO模型，针对红外人员及遮挡目标进行优化 \\

        RDK部署
        & 将模型转换并部署至RDK平台
        & 将模型转换并部署至RDK平台 \\

        部署测试
        & 测试检测精度、推理速度、延迟及资源占用
        & 测试检测精度、推理速度、延迟及资源占用 \\

        多路扩展
        & 使用多段视频模拟多路监控输入，测试多路推理性能
        & 使用多段红外视频模拟多路输入，测试并发推理性能 \\

        应用价值
        & 面向湿地生态保护和候鸟监测
        & 面向灾害环境搜救，辅助人员目标发现 \\

        数据获取难度
        & 中等，需要获得合法可用的鸟类图像或视频素材
        & 较低，POP数据集已公开，可直接获取 \\

        硬件需求
        & 不需要特殊硬件
        & 不需要无人机或红外摄像设备 \\

        技术难度
        & 中等，重点为多类别及小目标检测
        & 中等，重点为红外、遮挡及小目标检测 \\

        论文特色
        & 国内生态场景 + 候鸟检测 + 边缘AI部署
        & 红外成像 + 无人机搜救 + 遮挡检测 + 边缘AI部署 \\

        主要风险
        & 数据获取和数据规模可能影响实验效果
        & 红外数据与实际应用场景存在一定域差异 \\

        预期成果
        & 完成候鸟检测模型训练、优化及RDK部署
        & 完成人员红外检测模型训练、优化及RDK部署 \\

        \hline
    \end{tabularx}
\end{table}